\NativeSlot{50}{156}{%
% p076_s01

\ExamQuestion{4.}{7}%
学校对四年级的学生进行绘画能力抽查。三位老师的评分如下（单位：分）。\par




\par\noindent\makebox[\linewidth][c]{\begin{minipage}{.67\linewidth}\centering
\par\smallskip\Needspace{4\baselineskip}\begingroup\small\setlength{\tabcolsep}{3pt}
\centering\textbf{四（1）班学生绘画能力抽查表（单位：分）}\par\smallskip
\noindent\begin{tabularx}{\linewidth}{|*{7}{>{\centering\arraybackslash}X|}}\hline
 & 1号 & 2号 & 3号 & 4号 & 5号 & 6号\\\hline
王老师 & 90 & 88 & 95 & 65 & 76 & 86\\\hline
李老师 & 94 & 90 & 96 &  & 80 & 88\\\hline
刘老师 & 95 & 86 & 94 & 66 & 75 & 93\\\hline
平均得分 & 93 & 88 &  & 67 & 77 & 89\\\hline
\end{tabularx}\par
\endgroup\smallskip
\par\smallskip\Needspace{4\baselineskip}\begingroup\small\setlength{\tabcolsep}{3pt}
\centering\textbf{四（2）班学生绘画能力抽查表（单位：分）}\par\smallskip
\noindent\begin{tabularx}{\linewidth}{|*{7}{>{\centering\arraybackslash}X|}}\hline
 & 1号 & 2号 & 3号 & 4号 & 5号 & 6号\\\hline
王老师 & 83 & 93 & 93 & 83 & 88 & 77\\\hline
李老师 & 88 & 92 & 95 & 86 &  & 80\\\hline
刘老师 & 93 & 91 &  & 83 & 92 & 74\\\hline
平均得分 & 88 & 92 & 95 &  & 90 & 77\\\hline
\end{tabularx}\par
\endgroup\smallskip
\end{minipage}}\par
\noindent （1）将上面的表格填写完整。（10分）\par
}
\NativeSlot{403}{85}{%
% p076_s02

\ExamQuestion{（2）}{7}%
每个学生的平均得分若92～100分为优，80～91分为良，68～79分为达标，68分以下为待达标，将抽查成绩整理在下面的统计表中。（单位：人。）（8分）\par





\par\smallskip\Needspace{4\baselineskip}\begingroup\small\setlength{\tabcolsep}{3pt}
\noindent\textbf{统计表（单位：人）}\par\smallskip
\noindent\begin{tabularx}{\linewidth}{|*{5}{>{\centering\arraybackslash}X|}}\hline
班级 & 优 & 良 & 达标 & 待达标\\\hline
四（1）班 &  &  &  & \\\hline
四（2）班 &  &  &  & \\\hline
\end{tabularx}\par
\endgroup\smallskip
}
\NativeSlot{553}{37}{%
% p076_s03

\ExamQuestion{（3）}{7}%
结合统计表，分析一下两班学生绘画能力的情况。（5分）\par
}
